\section{Discussion}

All three primary candidates had higher mean correlations than the baseline,
but every confidence interval included zero. The layer-$-2$ and rich-statistics
heads improved on every seed. Their consistency across training runs did not
resolve uncertainty about the external participants.

The sign-flip tests and bootstrap intervals describe different sources of
variation. The former compare training runs on the fixed MIPDB cohort; the
latter also change which participants enter the evaluation. Small seed-level
p-values can therefore coexist with wide intervals when the estimated
advantage depends on the people being evaluated. Ten seeds provide ten training
runs, not ten independently recruited external cohorts.

\paragraph{The metric changes the preferred head.}
The rich-statistics head had the highest mean Pearson correlation
(\RichResidualPearson{}), whereas the multi-query head had the lowest MAE
(\MultiQueryMAE{} years). Pearson measures association, not agreement in years.
All four heads had negative external $R^2$, and their calibration slopes
exceeded one. Thus a moderate correlation did not imply accurate or
well-calibrated age estimates. These descriptive differences are useful for
understanding the heads, but they were not substituted for the prespecified
endpoint.

\paragraph{More training data did not yield a larger observed advantage.}
In the training-size analysis, the richer heads' mean advantages on MIPDB were
smaller at 800 training subjects than at 200. The endpoint confidence intervals
included zero, so this pattern alone does not establish that the advantage
shrinks with more data. Additional training data may improve the baseline and
candidate together; transfer between cohorts may also limit either predictor.
The experiment does not isolate these explanations. Choosing a new training
size from its MIPDB score would require a fresh external test.

The layer-wise analysis likewise found higher point estimates before the final
layer, but no interval excluded zero. The separate ds006780 transfer produced
a different training-size pattern: the rich head was below the baseline at
200 subjects and above it at 800. At 800, all ten differences were positive,
but the interval again included zero and exceeded the planned width. This
second cohort extends the comparison to another acquisition setting without
establishing a reliable advantage.

\paragraph{Implications.}
A frozen-encoder score describes an encoder used with a particular head,
training set, and metric. Reporting these choices, absolute prediction errors,
and participant-aware intervals makes comparisons easier to interpret. Our
primary, training-size, and layer-wise MIPDB analyses all reuse the same
participants; only the ds006780 extension adds a separate cohort. A larger
external evaluation with a protocol fixed in advance would test whether the
observed gains persist. Evaluating another encoder would address a different
question: how specific these findings are to REVE.
